% Reusable proof text, independently derived. Suggested bibliography keys:
% Nomura1994 (correct DOI 10.1007/BF02108300), Hatcher2002,
% Baez2002 (DOI 10.1090/S0273-0979-01-00934-X), FK1994.
% The real affine-kernel lemma below replaces the longer determinantal proof.

\subsection{Saturation and support orbits}
For a simple Euclidean Jordan algebra $V$ of rank $r$ and Peirce
constant $a$, write $\mathcal I_p(V)$ for its manifold of rank-$p$
idempotents. Its dimension is $ap(r-p)$ and its tangent space at $e$
is $V(e,1/2)$. These are classical facts about idempotent manifolds
\cite{FK1994,Nomura1994}. We use them to identify the additional restriction
imposed by globally smooth factor maps.

\begin{theorem}[Global saturation]\label{thm:global-saturation}
Let $C\subset\R^{n+1}$, $n\geq1$, contain the origin in its interior,
and suppose $P=\partial C$ and $D=\partial C^\circ$ are strictly convex
$C^1$ hypersurfaces. Suppose the full slack has a finite globally labelled
factorization
\[
 1-\ip{x}{z}=\sum_i\ip{X_i(x)}{Y_i(z)},\qquad (x,z)\in P\times D,
\]
where $X_i:P\to K_i$ and $Y_i:D\to K_i$ are $C^1$ and $K_i$ is a simple
symmetric cone with rank $r_i$ and Peirce constant $a_i$. Put
$c_i=a_i\lfloor r_i^2/4\rfloor$, with $c_i=0$ for rays.
If $\sum_i c_i=n$, each positive-capacity factor has constant balanced
complementary contact ranks $p_i,r_i-p_i$, and the joint support map
\[
 P\longrightarrow\prod_{i:c_i>0}\mathcal I_{p_i}(V_i),
 \qquad x\longmapsto(\operatorname{supp}X_i(x))_i
\]
is a finite covering. There is exactly one positive-capacity factor.
It is the Lorentz cone $Q_{n+2}$, except that $n=2$ also permits
$\mathbb S_+^3$ at the level of the support-orbit topology.
\end{theorem}
\begin{proof}
Let $z$ be the unique polar contact of $x$. The contact correspondence is
a homeomorphism; no differentiability of that correspondence is needed.
The tangent spaces are $T_xP=z^\perp$ and $T_zD=x^\perp$.
Their Euclidean pairing is nondegenerate: its left kernel consists of
vectors in $z^\perp\cap\operatorname{span}\{x\}$, which is zero because
$\ip{x}{z}=1$. Mixed differentiation in independent boundary charts gives
\[
 \ip{u}{v}=\sum_i H_i(u,v),\qquad
 H_i(u,v)=-\ip{dX_i(x)[u]}{dY_i(z)[v]}.
\]
The cross-Peirce argument of Lemma~\ref{lem:peirce-channel}, which for its
rank bound uses only the two independent derivative images, gives
\[
 \operatorname{rank}H_i\leq a_ip_iq_i\leq c_i,
 \quad p_i=\jrank X_i(x),\quad q_i=\jrank Y_i(z).
\]
Nondegeneracy and $\sum_i c_i=n$ force equality throughout. Thus
$p_i+q_i=r_i$, the ranks are balanced, and $\operatorname{rank}H_i=c_i$.
Ranks are lower semicontinuous along the contact graph. Since the two
ranks sum to $r_i$, each is locally constant, hence constant by
connectedness.

Let $e_i(x)=\operatorname{supp}X_i(x)$. This is $C^1$: locally the zero
spectral cluster is separated from the positive spectrum, and a smooth
scalar cutoff equal to zero on the former and one on the latter gives
$e_i$ by spectral functional calculus. Differentiating
$e_i\circ e_i=e_i$ and $X_i\circ e_i=X_i$ shows that
\[
 de_i[u]\in V_i(e_i,1/2),\qquad
 L_{X_i}de_i[u]=\tfrac12\operatorname{proj}_{V_i(e_i,1/2)}dX_i[u].
\]
Here $L_{X_i}$ is invertible on the half-Peirce space: in a frame adapted
to the positive support it acts on a cross channel by $\lambda_j/2>0$.
Consequently $de_i[u]=0$ precisely when the cross-Peirce component of
$dX_i[u]$ vanishes. In particular
\[
 c_i=\operatorname{rank}H_i\leq\operatorname{rank}de_i\leq c_i.
\]
If every $de_i$ kills $u$, every $H_i(u,v)$ is zero, so $u=0$.
The joint support differential is therefore injective, and its domain and
codomain have the same dimension $n$. The joint map is a local
$C^1$ diffeomorphism. Its image is open and, by compactness, closed in the
connected target. It is surjective and proper, hence a finite covering.

We give the topological classification explicitly. For $n\geq2$, the
source $P\cong S^n$ is simply connected, so this is the universal cover of
the target. There can be no circle factor, since it would make the target
fundamental group infinite. The remaining factors have compact simply
connected universal covers. The universal cover of their product is the
product of those covers. If there were at least two positive-dimensional
factors, the mod-two top cohomology class of one factor, tensored with the
unit classes of the others, would give nonzero cohomology in a degree
strictly between zero and $n$. The K\"unneth formula contradicts the
cohomology of $S^n$. Thus there is one factor. For $n=1$, the same conclusion
follows directly from the sum of the positive integer dimensions.

The simple-algebra classification now leaves the following cases.
A spin factor has primitive-idempotent orbit $S^{m-2}$ and therefore gives
$Q_{n+2}$. For $H_r(\R)$ the orbit is the real Grassmannian
$\operatorname{Gr}_p(\R^{p+q})$, $p+q=r$. At $r=3$ its balanced orbit is
$\mathbb{RP}^2$, covered by $S^2$. At $r\geq4$, balanced ranks satisfy
$p,q\geq2$. The oriented Grassmannian is simply connected by the exact
sequence of
\[
 SO(p)\times SO(q)\longrightarrow SO(p+q)
 \longrightarrow\operatorname{Gr}_p^+(\R^{p+q}):
\]
the induced map on fundamental groups onto $\pi_1(SO(p+q))=\mathbb Z_2$
is surjective. Its kernel is nonzero: use the diagonal element when
$p,q\geq3$, and an even nonzero winding in an $SO(2)$ factor otherwise.
Exactness therefore gives nonzero $\pi_2$ of the oriented Grassmannian.
Its dimension $pq\geq4$ precludes a sphere.

The complex Grassmannian is simply connected by the fibration with
isotropy $U(p)\times U(q)$ in $U(p+q)$. The map on $\pi_1$ is addition
$\mathbb Z^2\to\mathbb Z$, so its kernel gives nonzero $\pi_2$ of the
Grassmannian. For $r\geq3$ its dimension $2pq\geq4$ precludes a sphere.
For quaternionic Grassmannians, the groups $Sp(k)$ are simply connected
and $\pi_3(Sp(k))=\mathbb Z$, with standard inclusions preserving a
generator. These facts follow inductively from
$Sp(k-1)\to Sp(k)\to S^{4k-1}$ and $Sp(1)=S^3$.
The map on $\pi_3$ of the isotropy inclusion is again addition
$\mathbb Z^2\to\mathbb Z$, giving nonzero $\pi_4$ of the Grassmannian.
Its dimension $4pq\geq8$ precludes a sphere. The classical-group and
Grassmannian fibrations used here are described in
\cite[Examples~4.53--4.55]{Hatcher2002}.
Finally the balanced idempotent orbit of $H_3(\OO)$ is the Cayley plane:
rank-one idempotents are its points and complementation identifies the
rank-two orbit with it~\cite[Section~3.4]{Baez2002}.
Its cells have dimensions $0,8,16$, so $H^8(-;\mathbb Z)=\mathbb Z$;
see~\cite[Example~4.47]{Hatcher2002}. It is not a sphere.
Order-two real, complex, and quaternionic Hermitian algebras are already
spin factors. This exhausts the classification.
\end{proof}

\subsection{Eliminating the real order-three exception}
The apparent exceptional possibility is ruled out by the full ball slack.
The following elementary observation avoids a classification of linear
determinantal representations.

\begin{lemma}[Injective kernels of an affine sphere pencil]
\label{lem:affine-kernel-injective}
Let $s\geq2$ and $X(x)=C_0+\sum_{j=1}^s x_jC_j$ be a real symmetric
matrix pencil, positive semidefinite of nullity one for every
$x\in S^{s-1}$. If its kernels do not share a nonzero vector, its
kernel-line map is injective.
\end{lemma}
\begin{proof}
Suppose distinct $x,y$ have the same kernel line, containing $k\ne0$.
The affine scalar function $f(v)=k^TX(v)k=a+b\cdot v$ is nonnegative on
the unit sphere and vanishes at $x$ and $y$. Nonnegativity says
$a\geq\|b\|$. If $b\ne0$, a zero is possible only when
$a=\|b\|$ and $v=-b/\|b\|$, giving only one zero. Thus $a=b=0$.
For every $v$, positivity and $k^TX(v)k=0$ imply $X(v)k=0$, contradicting
the absence of a common nonzero kernel vector.
\end{proof}

\begin{lemma}[Finite affine selection on a sphere]
\label{lem:finite-affine-selection}
Let $n\geq2$ and let $F:S^n\to W$ be $C^1$, with $W$ finite dimensional.
Suppose a dense open subset is the union of finitely many nonempty open
sets $U_i$, and $F|_{U_i}=L_i|_{U_i}$ for ambient affine maps
$L_i:\R^{n+1}\to W$. Then $F$ is the restriction of one affine map.
\end{lemma}
\begin{proof}
Merge sets with identical $L_i$ and put $E_i=\overline{U_i}$. Continuity
of values and tangent derivatives gives $j^1F=j^1(L_i|_{S^n})$ on $E_i$.
If distinct affine maps have the same tangent first jet at $z\in S^n$,
their difference is $M(z\cdot x-1)$ for some nonzero $M\in W$.
Such a difference has matching first jet at no other sphere point:
matching derivatives force the other normal line to be $\R z$, and
matching values exclude $-z$. Thus each intersection $E_i\cap E_j$ lies
in at most one point. Remove the resulting finite set $Z$. The connected
space $S^n\setminus Z$ is partitioned into the disjoint relatively closed
sets $E_i\setminus Z$; each is also relatively open because the partition
is finite. Only one is nonempty. But the closure of a nonempty open $U_i$
cannot be finite, so there was only one distinct affine map. Density
finishes the proof.
\end{proof}

\begin{proposition}[No real order-three saturation for a ball]
\label{prop:no-real-three-saturation}
There is no factorization
\[
 1-x\cdot y=\operatorname{tr}(X(x)Y(y))+\sum_{j=1}^q a_j(x)b_j(y),
 \qquad x,y\in S^2,
\]
with $X,Y:S^2\to\mathbb S_+^3$ of class $C^1$, continuous nonnegative
$a_j,b_j$, finite $q$, constant complementary ranks $2,1$, and a
rank-one support map that is a local diffeomorphism onto $\mathbb{RP}^2$.
\end{proposition}
\begin{proof}
Exchange $X,Y$ if necessary so that $Y$ has rank one. Every diagonal
summand is zero, hence $a_j(x)b_j(x)=0$. For each of the finitely many
continuous nonnegative functions $f=a_j,b_j$, the set
$\{f>0\}\cup\operatorname{int}\{f=0\}$ is dense open. Intersect these
sets and partition the result by their finitely many positivity/zero
patterns. On each nonempty open pattern set $U$, for each $j$ at least
one of $a_j|_U,b_j|_U$ is identically zero. Therefore all ray terms vanish
for every pair $(x,y)\in U\times U$.

The rank-one support map of $Y$ sends $U$ onto a projective open set.
Consequently the matrices $Y(y)$, $y\in U$, span $\mathbb S^3$: an
annihilating symmetric matrix would have a quadratic form zero on an open
set of projective directions, and hence zero identically. Choose six
$y_\ell\in U$ whose matrices are a basis. The equations
\[
 \operatorname{tr}(X(x)Y(y_\ell))=1-x\cdot y_\ell,
 \qquad x\in U,
\]
force $X|_U$ to be an ambient affine pencil. There are only finitely many
pattern sets, so Lemma~\ref{lem:finite-affine-selection} makes $X$ one
affine pencil on the whole sphere.

Complementarity identifies the kernel line of $X(x)$ with the support
line of $Y(x)$. This map is a finite covering $S^2\to\mathbb{RP}^2$,
hence the universal double cover. Its image contains all lines, so its
kernels have no common nonzero vector. But
Lemma~\ref{lem:affine-kernel-injective} makes the map injective,
contradicting its two-point fibers.
\end{proof}

\begin{corollary}[Ball saturation classification]
\label{cor:ball-global-saturation}
A finite globally bi-$C^1$ symmetric-cone factorization of the full slack
of $\B_2^s$, $s\geq2$, satisfying
$\sum_i a_i\lfloor r_i^2/4\rfloor=s-1$ has exactly one
positive-capacity factor, namely $Q_{s+1}$. Remaining factors, if any,
are rays.
\end{corollary}
\begin{proof}
Apply Theorem~\ref{thm:global-saturation} and exclude its sole additional
case by Proposition~\ref{prop:no-real-three-saturation}.
\end{proof}

\subsection{A sharp Hermitian contact-rank bound}
The same support argument yields a different saturation theorem for
selected certificate ranks. Let $\mathbb F\in\{\R,\C,\HH\}$,
$\delta=\dim_\R\mathbb F$, $R\geq2$, and $B=\delta(R-1)$. Define
\[
 \kappa_{\mathbb F,R}(s)=
 \begin{cases}
 1,&R=2\text{ and }s=\delta+1,\\
 \lceil s/B\rceil,&\text{otherwise}.
 \end{cases}
\]
\begin{proposition}[Hermitian contact rank]\label{prop:hermitian-contact-rank}
Suppose the full slack of $\B_2^s$, $s\geq2$, has a finite globally
bi-$C^1$ factorization with factors in $H_+^{r_i}(\mathbb F)$,
$r_i\leq R$. Some $v\in S^{s-1}$ satisfies
$\sum_i\operatorname{rank}_{\mathbb F}Y_i(v)
\geq\kappa_{\mathbb F,R}(s)$.
\end{proposition}
\begin{proof}
Write $Q(v)=\sum_i q_i(v)$, $q_i(v)=\operatorname{rank}_{\mathbb F}Y_i(v)$.
The mixed-channel bound at every contact gives
\[
 s-1\leq\sum_i\delta p_iq_i
 \leq\sum_i\delta(R-q_i)q_i\leq BQ(v).
\]
If $B$ does not divide $s-1$, the ceiling of $(s-1)/B$ equals
$\lceil s/B\rceil$, proving the result. Otherwise write $s-1=qB$ and
suppose $Q(v)\leq q$ everywhere. Then $Q\equiv q$. Lower
semicontinuity makes all individual dual ranks constant. Equality in the
three bounds forces every nonzero dual block to have rank one, order
$R$, and primal rank $R-1$; every other dual block vanishes identically.
There are exactly $q$ nonzero dual blocks. The support differential and
joint-injectivity argument from Theorem~\ref{thm:global-saturation} gives
a finite covering
\[
 S^{qB}\longrightarrow(\mathbb F P^{R-1})^q.
\]
If $q>1$, a circle factor or the product cohomology argument in that proof
rules this out. If $q=1$, complex and quaternionic projective spaces have
nonzero cohomology in degrees two and four respectively, and can be
spheres only at $R=2$. For the real case with $R\geq3$, the covering is
$S^{R-1}\to\mathbb{RP}^{R-1}$ and has two-point fibers.
Only one matrix factor remains in the full slack identity. Its rank-one
sheet spans $\mathbb S^R$, by the quadratic-form argument in
Proposition~\ref{prop:no-real-three-saturation}. Pairing against a fixed
basis of that sheet therefore makes the rank-$(R-1)$ primal sheet one
affine pencil globally. Its kernel lines cover all of
$\mathbb{RP}^{R-1}$, so there is no common nonzero kernel vector.
Lemma~\ref{lem:affine-kernel-injective} contradicts the double cover.
Thus the sole surviving case is $q=1,R=2,s=\delta+1$. In all other cases
some $Q(v)\geq q+1=\lceil s/B\rceil$.
\end{proof}

% The product-ball sequential compression proof belongs with the parent
% author's products subsection: apply Proposition above after each
% orthogonal compression and telescope the native-field range dimensions.
% The bound concerns these globally selected factors. It is not a lower
% bound on the minimum rank over a certificate fiber with no global C1
% selection. The no-selection invariant was treated earlier in the paper.
